\begin{abstract}
Carbon monoxide binds strongly to platinum and can deactivate Pt-based
catalysts. Ge incorporation can weaken this interaction, but it also changes
the geometry and bonding pattern of small Pt--Ge clusters. We use ten-atom
\ce{Pt10}, \ce{Pt6Ge4}, \ce{Pt5Ge5}, and
\ce{Pt4Ge6} models to examine how these structural and electronic effects act
together.

The main comparison is based on crossed-parent models. In one route, Ge atoms
are introduced into a common \ce{Pt10} framework; in the reverse route, a
\ce{Pt5Ge5} framework is converted into a pure \ce{Pt10} derivative. Within
the common \ce{Pt10}-parent series, the mean binding energy becomes less
negative as the Ge content increases. AdNDP shows a change from delocalized
multicenter bonding in pure \ce{Pt10} to more localized Pt--Ge bonding in the
Ge-containing models. {Calculations with \ce{CO} placed at every
independent Pt site show that the affinity of an individual site depends on the
local environment, parent geometry, and cluster reorganization. Calculations
initiated from previously reported global-minimum structures further show that
the adsorption landscape changes with the parent isomer and spin state.
Overall, increasing the Ge content weakens mean \ce{CO} adsorption within the
crossed-parent series. Individual binding energies vary with the local
environment and with the structural response of the cluster after \ce{CO}
adsorption.}
\end{abstract}
